\documentclass[10pt,a4paper]{amsart}
\usepackage{amsmath,amssymb,amsthm,graphicx}
\usepackage{verbatim}
\newenvironment{reference}{\comment}{\endcomment}
%\newenvironment{reference}{}{}
\newenvironment{slogan}{\comment}{\endcomment}
\newenvironment{history}{\comment}{\endcomment}

% For commutative diagrams we use Xy-pic
\usepackage[all]{xy}

% We use 2cell for 2-commutative diagrams.
\xyoption{2cell}
% Theorem environments.
%
\theoremstyle{plain}
\newtheorem{theorem}[subsection]{Theorem}
\newtheorem{proposition}[subsection]{Proposition}
\newtheorem{lemma}[subsection]{Lemma}

\theoremstyle{definition}
\newtheorem{definition}[subsection]{Definition}
\newtheorem{example}[subsection]{Example}
\newtheorem{exercise}[subsection]{Exercise}
\newtheorem{situation}[subsection]{Situation}

\theoremstyle{remark}
\newtheorem{remark}[subsection]{Remark}
\newtheorem{remarks}[subsection]{Remarks}

\numberwithin{equation}{subsection}

% Macros
%
\def\lim{\mathop{\mathrm{lim}}\nolimits}
\def\colim{\mathop{\mathrm{colim}}\nolimits}
\def\Spec{\mathop{\mathrm{Spec}}}
\def\Hom{\mathop{\mathrm{Hom}}\nolimits}
\def\Ext{\mathop{\mathrm{Ext}}\nolimits}
\def\SheafHom{\mathop{\mathcal{H}\!\mathit{om}}\nolimits}
\def\SheafExt{\mathop{\mathcal{E}\!\mathit{xt}}\nolimits}
\def\Sch{\mathit{Sch}}
\def\Mor{\mathop{\mathrm{Mor}}\nolimits}
\def\Ob{\mathop{\mathrm{Ob}}\nolimits}
\def\Sh{\mathop{\mathit{Sh}}\nolimits}
\def\NL{\mathop{N\!L}\nolimits}
\def\CH{\mathop{\mathrm{CH}}\nolimits}
\def\proetale{{pro\text{-}\acute{e}tale}}
\def\etale{{\acute{e}tale}}
\def\QCoh{\mathit{QCoh}}
\def\Ker{\mathop{\mathrm{Ker}}}
\def\Im{\mathop{\mathrm{Im}}}
\def\Coker{\mathop{\mathrm{Coker}}}
\def\Coim{\mathop{\mathrm{Coim}}}

% Boxtimes
%
\DeclareMathSymbol{\boxtimes}{\mathbin}{AMSa}{"02}

%
% Macros for moduli stacks/spaces
%
\def\QCohstack{\mathcal{QC}\!\mathit{oh}}
\def\Cohstack{\mathcal{C}\!\mathit{oh}}
\def\Spacesstack{\mathcal{S}\!\mathit{paces}}
\def\Quotfunctor{\mathrm{Quot}}
\def\Hilbfunctor{\mathrm{Hilb}}
\def\Curvesstack{\mathcal{C}\!\mathit{urves}}
\def\Polarizedstack{\mathcal{P}\!\mathit{olarized}}
\def\Complexesstack{\mathcal{C}\!\mathit{omplexes}}
% \Pic is the operator that assigns to X its picard group, usage \Pic(X)
% \Picardstack_{X/B} denotes the Picard stack of X over B
% \Picardfunctor_{X/B} denotes the Picard functor of X over B
\def\Pic{\mathop{\mathrm{Pic}}\nolimits}
\def\Picardstack{\mathcal{P}\!\mathit{ic}}
\def\Picardfunctor{\mathrm{Pic}}
\def\Deformationcategory{\mathcal{D}\!\mathit{ef}}
\usepackage{cleveref}
\crefname{originalreference}{original reference}{original references}
\Crefname{originalreference}{Original reference}{Original references}
\makeatletter
\newlabel{lemma-get-w-local-algebraic-residue-field-extensions}{{097X}{source}{}{sourcepointer-1}{}}
\newlabel{lemma-get-w-local-algebraic-residue-field-extensions@cref}{{[originalreference][0][]097X}{[1][0][]source}{}{}{}}
\makeatother
\crefname{originallemma}{original lemma}{original lemmas}
\Crefname{originallemma}{Original lemma}{Original lemmas}
\makeatletter
\newlabel{more-algebra-lemma-weakly-etale-permanence}{{092L}{source}{}{sourcepointer-2}{}}
\newlabel{more-algebra-lemma-weakly-etale-permanence@cref}{{[originallemma][0][]092L}{[1][0][]source}{}{}{}}
\makeatother
\crefname{originallemma}{original lemma}{original lemmas}
\Crefname{originallemma}{Original lemma}{Original lemmas}
\makeatletter
\newlabel{more-algebra-lemma-weakly-etale-residue-field-extensions}{{092R}{source}{}{sourcepointer-3}{}}
\newlabel{more-algebra-lemma-weakly-etale-residue-field-extensions@cref}{{[originallemma][0][]092R}{[1][0][]source}{}{}{}}
\makeatother
\crefname{originaltheorem}{original theorem}{original theorems}
\Crefname{originaltheorem}{Original theorem}{Original theorems}
\makeatletter
\newlabel{more-algebra-theorem-olivier}{{092Z}{source}{}{sourcepointer-4}{}}
\newlabel{more-algebra-theorem-olivier@cref}{{[originaltheorem][0][]092Z}{[1][0][]source}{}{}{}}
\makeatother
\crefname{originalreference}{original reference}{original references}
\Crefname{originalreference}{Original reference}{Original references}
\makeatletter
\newlabel{proposition-maps-wich-identify-local-rings}{{097G}{source}{}{sourcepointer-5}{}}
\newlabel{proposition-maps-wich-identify-local-rings@cref}{{[originalreference][0][]097G}{[1][0][]source}{}{}{}}
\makeatother
\title{A weakly etale ring map becomes ind-etale after a faithfully flat ind-etale extension}
\author{The Stacks Project Authors}
\date{stacks-project:proetale.tex:local-snapshot-20261012; GNU FDL 1.2 or later}
\begin{document}
\maketitle
\noindent\textit{Editorial note.} All delivered mathematical text is the author's own, copied from the pinned Stacks Project chapter file proetale.tex; no statement or proof step was rewritten, reordered or removed.\par\medskip
\noindent\textit{Editorial note.} The counted proof is the author's complete argument: the passage to the w-local diagram with strictly henselian local rings, the proof that the map between the middle and the top identifies local rings through rigidity of weakly etale maps between strictly henselian local rings, and the closing application of the local isomorphism proposition.\par\medskip
\noindent\textit{Editorial note.} The label proposition-weakly-etale names the counted statement, and the proof environment immediately following it in the source is delivered verbatim as its proof.\par\medskip
\noindent\textit{Editorial note.} The counted statement is carried by the author's own proposition environment in the chapter source.\par\medskip
\noindent\textit{Editorial note.} Pre-open gate: python operations/agent-authoring/tools/precard\_gate.py --source proetale.tex --label proposition-weakly-etale --start 1801 --end 1834 --local-id 3595 returned RESULT=OK (L3595 free, occupancy FREE, public interval clear). Candidate taken from sources/stacks-proetale/candidates-halfA.json.\par\medskip
\noindent\textit{Editorial note.} The statement and proof are copied from the pinned chapter file; the cited inputs are kept as external pointers with their canonical Stacks tags.\par\medskip
\noindent\textit{Allowed standard prerequisites (editorial).}\begin{itemize}
\item Weakly etale maps induce algebraic residue field extensions. More on Algebra, the lemma that a weakly etale ring map induces separable algebraic extensions of residue fields. Supplies the hypothesis of the residue field lemma applied to A -> B.
\item The diagram with w-local spectra and strictly henselian local rings. Same chapter, the lemma producing a diagram A -> C -> D over A -> B with w-local spectra, closed point conditions and strictly henselian local rings. Supplies the diagram whose properties are used throughout the proof.
\item Rigidity of weakly etale maps between strictly henselian local rings. More on Algebra, the theorem that a weakly etale map of strictly henselian local rings is an isomorphism, and the lemma on permanence of the weakly etale property. Supply the isomorphism C\_m' -> D\_m at maximal ideals and the weak etaleness of C -> D.
\item Faithfully flat ind-Zariski maps making a local isomorphism ind-Zariski. Same chapter, the proposition that a ring map which identifies local rings can be made ind-Zariski after a faithfully flat ind-Zariski extension. Supplies the map D -> D' with C -> D' ind-Zariski, which finishes the proof.
\end{itemize}\medskip
\ifdefined\hypertarget\hypertarget{sourcepointer-1}{}\fi
\noindent\textit{Source reference.} Original reference 097X: sources/stacks-proetale/pinned-original-source.tex, source lines 1723--1723. Prerequisite pointer invoked at proof line(s) 1813; it names a standard fact of the same chapter that the counted argument uses, and every step of the counted proof built on it is retained verbatim in the delivered body.\par\medskip
\ifdefined\hypertarget\hypertarget{sourcepointer-2}{}\fi
\noindent\textit{Source reference.} Original lemma 092L: sources/stacks-chow/tags.txt, source lines 11484--11484. Prerequisite pointer invoked at proof line(s) 1823; it names a standard fact of a sibling chapter that the counted argument uses, and every step of the counted proof built on it is retained verbatim in the delivered body.\par\medskip
\ifdefined\hypertarget\hypertarget{sourcepointer-3}{}\fi
\noindent\textit{Source reference.} Original lemma 092R: sources/stacks-chow/tags.txt, source lines 11489--11489. Prerequisite pointer invoked at proof line(s) 1811; it names a standard fact of a sibling chapter that the counted argument uses, and every step of the counted proof built on it is retained verbatim in the delivered body.\par\medskip
\ifdefined\hypertarget\hypertarget{sourcepointer-4}{}\fi
\noindent\textit{Source reference.} Original theorem 092Z: sources/stacks-chow/tags.txt, source lines 11497--11497. Prerequisite pointer invoked at proof line(s) 1826; it names a standard fact of a sibling chapter that the counted argument uses, and every step of the counted proof built on it is retained verbatim in the delivered body.\par\medskip
\ifdefined\hypertarget\hypertarget{sourcepointer-5}{}\fi
\noindent\textit{Source reference.} Original reference 097G: sources/stacks-proetale/pinned-original-source.tex, source lines 1228--1228. Prerequisite pointer invoked at proof line(s) 1830; it names a standard fact of the same chapter that the counted argument uses, and every step of the counted proof built on it is retained verbatim in the delivered body.\par\medskip
\begin{proposition}
\label{proposition-weakly-etale}
Let $A \to B$ be a weakly \'etale ring map.
Then there exists a faithfully flat, ind-\'etale ring map
$B \to B'$ such that $A \to B'$ is ind-\'etale.
\end{proposition}

\begin{proof}
The ring map $A \to B$ induces (separable) algebraic extensions of
residue fields, see More on Algebra, Lemma
\ref{more-algebra-lemma-weakly-etale-residue-field-extensions}.
Thus we may apply
Lemma \ref{lemma-get-w-local-algebraic-residue-field-extensions}
and choose a diagram
$$
\xymatrix{
B \ar[r] & D \\
A \ar[r] \ar[u] & C \ar[u]
}
$$
with the properties as listed in the lemma. Note that $C \to D$
is weakly \'etale by
More on Algebra, Lemma \ref{more-algebra-lemma-weakly-etale-permanence}.
Pick a maximal ideal $\mathfrak m \subset D$. By construction
this lies over a maximal ideal $\mathfrak m' \subset C$.
By More on Algebra, Theorem \ref{more-algebra-theorem-olivier}
the ring map $C_{\mathfrak m'} \to D_\mathfrak m$ is an isomorphism.
As every point of $\Spec(C)$ specializes to a closed point we conclude that
$C \to D$ identifies local rings.
Thus Proposition \ref{proposition-maps-wich-identify-local-rings}
applies to the ring map $C \to D$. Pick $D \to D'$ faithfully flat
and ind-Zariski such that $C \to D'$ is ind-Zariski. Then
$B \to D'$ is a solution to the problem posed in the proposition.
\end{proof}

\end{document}
